% =====================================================================
% Section 4: Experiments.
% Table bodies are generated by tables/gen_tables.py from result JSONs
% (hand-merged where marked [length-edit]); this file only \input's them
% and discusses labels. Numbers quoted in prose are tagged \RLOCAL and
% come from result files read at writing time; small-n pilots are
% explicitly flagged with their n and carry no significance claims.
% Notation: q = Hutchinson probes, N = ensemble/particle count,
% T = trajectory steps, R_a = acquisition rounds;
% NFE = T + 2q x R_a (selection-only, Eq. eq:nfe), 2T + 2q R_a for duo,
% N*T for particle ensembles.
% =====================================================================
\section{Experiments}\label{sec:exp}

We instantiate \pvas{} on block-pixel acquisition on natural images at
two scales --- CIFAR-10 (Secs.~\ref{sec:exp_sr}--\ref{sec:exp_inpaint})
and CelebA-HQ $256^2$ (Sec.~\ref{sec:exp_celeba}) --- and on multi-coil
MRI $k$-space line acquisition (Sec.~\ref{sec:exp_mri}), following
Sec.~\ref{sec:method_inst} (Alg.~S1, supplement). The goal is twofold:
measure reconstruction quality at a fixed measurement budget, and price
the uncertainty signal itself. Ensemble-based selection is therefore not
only cited but reimplemented (Sec.~\ref{sec:method_active}) and run
head-to-head in the same pipeline: same prior, same schedule, same
rounds, same budget.

\subsection{Protocol, statistics, and provenance}

\paragraph{Matched budget and matched \nfe.} All policies in a table
start from the same initial observations and end with the same number
of acquired atoms. They differ only in \emph{which} atoms and \emph{when}
(static: full budget up-front; active: spent mid-trajectory). New
measurements are read off the held-out ground truth plus fresh noise ---
the standard retrospective protocol of active-MRI evaluation
\cite{zhang2019active,bakker2020greedy}; the ground truth enters the
sampler in no other way. Every table carries an \nfe{} column counting
\emph{all} denoiser evaluations per image or slice (Eq.~\ref{eq:nfe})
--- our probes and every particle of ensemble-based selection: no
uncertainty signal is treated as free.

\paragraph{Statistics and provenance.} Within a configuration, all
methods see identical images, masks, noise realizations, and seeds, so
comparisons are paired; we report mean$\pm$std of \psnrm{}/\ssimm{}
\cite{wang2004ssim} and one-sided Wilcoxon signed-rank tests on paired
per-image / per-slice \psnrm{} deltas; $n$ is in each caption, seeds in
the released result files. \RLOCAL{} marks results reproduced with the
code in the supplement on a single Apple M4 Pro (MPS/CPU); \CITED{}
marks numbers from the cited paper, not re-run by us; ``\PENDING''
marks a cell whose run is fully specified but whose result file was not
yet written when this text was frozen. Small-$n$ pilots are always
quoted with their $n$ and never carry significance claims. No projected
or estimated numbers appear anywhere in the paper. Code and result
files will be released.

\subsection{Guided super-resolution on CIFAR-10}\label{sec:exp_sr}

\paragraph{Setup.} Following D2 (Sec.~\ref{sec:method_inst}):
$32{\times}32$ CIFAR-10, $T{=}100$ deterministic DDIM steps, $64$
blocks of $4{\times}4$ pixels as atoms; the $4{\times}$ box-downsampled
$8{\times}8$ image is \emph{always} observed, and the budget buys
$L{=}12$ full-resolution blocks ($18.75\%$ of the image). Measurements
are noiseless ($\sigma_y{=}0$), so consistency is hard replacement plus
LR back-projection. Active policies spend the
budget in $R_a{=}3$ rounds at $45\%$, $60\%$, and $75\%$ of the reverse
trajectory; \pvas{} runs selection-only with $q{=}4$ probes and the
smoothed block-sum score ($\nfe{=}100+2qR_a{=}124$); $n{=}128$ fixed
test images (Table~\ref{tab:sr_main}).

\paragraph{Policies at matched budget.}
\emph{static} (all $12$ blocks up-front, uniformly at random);
\emph{random}/\emph{raster} (mid-trajectory, same 3-round schedule,
uniform-random / fixed-scan-order scores); \emph{pv} (ours);
\emph{oracle} (blocks ranked by true squared error of the current
$\hat{\xz}$ --- a retrospective upper bound); \emph{ensemble}
(the faithful ADS/AdaSense-style baseline of
Sec.~\ref{sec:method_active} \cite{nolan2025ads,elata2024adasense}:
$N{=}8$ parallel posterior trajectories, blocks ranked by
cross-particle variance, reconstruction is the particle mean, $NT$
\nfe); and \emph{ens\_random} ($N{=}8$ particles, \emph{random}
selection --- isolating posterior-mean averaging from selection).

\begin{figure*}[t]
  \centering
  \begin{minipage}[t]{0.40\textwidth}
    \centering
    \includegraphics[width=0.72\linewidth]{figs/fig_budget.pdf}
    \caption{Budget sweep on CIFAR-10 guided SR: \psnrm{} vs.\ acquired
    blocks $L\in\{6,12,18,24\}$ at matched budget; $n{=}64$ images per
    point ($n{=}128$ at the main setting $L{=}12$)\RLOCAL.}
    \label{fig:budget}
  \end{minipage}\hfill
  \begin{minipage}[t]{0.56\textwidth}
    \centering
    \includegraphics[width=0.72\linewidth]{figs/fig_calibration.pdf}
    \caption{Calibration against the oracle block-error ranking:
    per-image Spearman correlation and top-$12$ overlap vs.\ extra
    \nfe{} for \pvas{} ($q\in\{1,\dots,16\}$ probes) and trajectory
    ensembles ($N\in\{4,8,16\}$ particles), in the SR and from-scratch
    inpainting regimes.}
    \label{fig:calibration}
  \end{minipage}
\end{figure*}

\begin{table*}[t]
  \begin{minipage}[t]{0.492\textwidth}
    % Auto-generated by paper/pvas_v2/tables/gen_tables.py -- do not edit.
% input: outputs/poc_natural/cifar_sr_main/results.json [R-LOCAL]
% n = 128 images (from JSON config); seed = 0
% [length-edit] shell stripped: lives inside a paired table* in experiments.tex
  \centering\small
  \setlength{\tabcolsep}{4pt}
  \captionof{table}{Adaptive block super-resolution on CIFAR-10 at matched measurement budget, $L{=}12$ blocks of 4$\times$4 px, $T{=}100$ steps, $\sigma_y{=}0.0$; $n{=}128$ images\RLOCAL. Oracle and ensemble rows are references and excluded from bolding. $\Delta$ rows are paired per-image \psnrm{} differences with one-sided Wilcoxon $p$.}
  \label{tab:sr_main}
  \begin{tabular}{lccc}
    \toprule
    Acquisition policy & \psnrm{} (dB) & \ssimm{} & \nfe/image \\
    \midrule
    Static mask (random, up-front) & 19.79\,$\pm$\,2.50 & 0.651\,$\pm$\,0.100 & 100 \\
    Random (mid-trajectory) & 19.98\,$\pm$\,2.46 & 0.660\,$\pm$\,0.097 & 100 \\
    Raster (mid-trajectory) & 19.38\,$\pm$\,2.27 & 0.603\,$\pm$\,0.097 & 100 \\
    \pvas{} (ours) & \textbf{21.43\,$\pm$\,3.70} & \textbf{0.701\,$\pm$\,0.106} & 124 \\
    Oracle error (upper bound) & 22.95\,$\pm$\,3.80 & 0.753\,$\pm$\,0.102 & 100 \\
    Ensemble-$8$ variance (ref.) & 25.42\,$\pm$\,3.88 & 0.825\,$\pm$\,0.075 & 800 \\
    Ensemble-$8$, random select (ref.) & 24.13\,$\pm$\,2.69 & 0.819\,$\pm$\,0.058 & 800 \\
    \midrule
    \multicolumn{4}{l}{\emph{Paired $\Delta$\psnrm{} (\pvas{} $-$ baseline), dB}} \\
    $\Delta$ vs.\ static & $+1.63$\,$\pm$\,$1.84$ & $p{=}1.0{\times}10^{-21}$ & \\
    $\Delta$ vs.\ random & $+1.45$\,$\pm$\,$2.01$ & $p{=}7.1{\times}10^{-18}$ & \\
    $\Delta$ vs.\ raster & $+2.05$\,$\pm$\,$2.20$ & $p{=}6.9{\times}10^{-22}$ & \\
    $\Delta$ vs.\ ensemble & $-3.99$\,$\pm$\,$1.45$ & $p{=}1.000$ & \\
    $\Delta$ vs.\ oracle & $-1.52$\,$\pm$\,$1.04$ & $p{=}1.000$ & \\
    \bottomrule
  \end{tabular}

       % -> tab:sr_main
  \end{minipage}\hfill
  \begin{minipage}[t]{0.492\textwidth}
    % Auto-generated by paper/pvas_v2/tables/gen_tables.py -- do not edit.
% input: outputs/poc_natural/cifar_inpaint_main/results.json [R-LOCAL]
% n = 128 images (from JSON config); seed = 0
% [length-edit] shell stripped: lives inside a paired table* in experiments.tex
  \centering\small
  \setlength{\tabcolsep}{4pt}
  \captionof{table}{Adaptive block inpainting on CIFAR-10 at matched measurement budget, $L{=}12$ blocks of 4$\times$4 px, $T{=}100$ steps, $\sigma_y{=}0.0$; $n{=}128$ images\RLOCAL. Unlike the SR task, no low-resolution preview is observed: reconstruction proceeds from scratch, conditioned only on the initial random fraction (0.1) of blocks plus the acquired blocks. Oracle and ensemble rows are references and excluded from bolding. $\Delta$ rows are paired per-image \psnrm{} differences with one-sided Wilcoxon $p$.}
  \label{tab:inpaint}
  \begin{tabular}{lccc}
    \toprule
    Acquisition policy & \psnrm{} (dB) & \ssimm{} & \nfe/image \\
    \midrule
    Static mask (random, up-front) & \textbf{14.23\,$\pm$\,2.32} & \textbf{0.401\,$\pm$\,0.092} & 100 \\
    Random (mid-trajectory) & 12.71\,$\pm$\,2.25 & 0.308\,$\pm$\,0.084 & 100 \\
    \pvas{} (ours) & 12.69\,$\pm$\,2.63 & 0.356\,$\pm$\,0.098 & 124 \\
    Oracle error (upper bound) & 15.72\,$\pm$\,2.80 & 0.436\,$\pm$\,0.107 & 100 \\
    Ensemble-$8$ variance (ref.) & 16.91\,$\pm$\,2.58 & 0.453\,$\pm$\,0.080 & 800 \\
    Ensemble-$8$, random select (ref.) & 17.30\,$\pm$\,2.36 & 0.525\,$\pm$\,0.064 & 800 \\
    \midrule
    \multicolumn{4}{l}{\emph{Paired $\Delta$\psnrm{} (\pvas{} $-$ baseline), dB}} \\
    $\Delta$ vs.\ static & $-1.54$\,$\pm$\,$1.60$ & $p{=}1.000$ & \\
    $\Delta$ vs.\ random & $-0.02$\,$\pm$\,$0.89$ & $p{=}0.848$ & \\
    $\Delta$ vs.\ ensemble & $-4.23$\,$\pm$\,$2.30$ & $p{=}1.000$ & \\
    $\Delta$ vs.\ oracle & $-3.03$\,$\pm$\,$1.28$ & $p{=}1.000$ & \\
    \bottomrule
  \end{tabular}

  % -> tab:inpaint
  \end{minipage}
\end{table*}

\begin{table*}[t]
  \begin{minipage}[t]{0.492\textwidth}
    % Auto-generated by paper/pvas_v2/tables/gen_tables.py -- do not edit.
% input: outputs/poc_natural/cifar_sr_noisy_nogate/results.json [R-LOCAL]
% input: outputs/poc_natural/cifar_sr_noisy_gate/results.json [R-LOCAL]
% n = 64 images; sigma_y = 0.1 (from JSON config)
% [length-edit] shell stripped: lives inside a paired table* in experiments.tex
  \centering\small
  \setlength{\tabcolsep}{4pt}
  \captionof{table}{Gate$\times$acquisition ablation on noisy CIFAR-10 super-resolution ($\sigma_y{=}0.1$, $n{=}64$ images)\RLOCAL: the same PV map used as acquisition score, as Kalman gain $\gain_i = \PV_i/(\PV_i+\sigma_i^2)$ for data consistency, or both. Best entry per column in bold (reference rows excluded). $\Delta$ row: paired per-image \psnrm{} difference, one-sided Wilcoxon $p$.}
  \label{tab:ablation_gate}
  \begin{tabular}{llccc}
    \toprule
    Acquisition & DC gain & \psnrm{} (dB) & \ssimm{} & \nfe/image \\
    \midrule
    Static (up-front) & scalar & 18.78\,$\pm$\,1.68 & 0.564\,$\pm$\,0.091 & 100 \\
    Random (mid-traj.) & scalar & 19.08\,$\pm$\,1.89 & 0.570\,$\pm$\,0.084 & 100 \\
    \textsc{pv} (ours) & scalar & \textbf{20.42\,$\pm$\,2.45} & 0.650\,$\pm$\,0.098 & 124 \\
    \textsc{pv} (ours) & \textsc{pv}-gated & 20.40\,$\pm$\,2.32 & \textbf{0.652\,$\pm$\,0.093} & 276 \\
    \midrule
    \multicolumn{2}{l}{$\Delta$(gated $-$ scalar), \textsc{pv} policy} & $-0.02$\,$\pm$\,$0.79$ & $p{=}0.287$ & \\
    \bottomrule
  \end{tabular}

 % -> tab:ablation_gate
  \end{minipage}\hfill
  \begin{minipage}[t]{0.492\textwidth}
    % Hand-restructured: every ensemble row is now measured head-to-head in our
% pipeline (run_ensemble_active / run_ensemble_active_mc); the published
% multipliers (ADS Np=16, AdaSense s=256) are one \CITED sentence in the
% Sec. cost prose, no longer table rows.
% input: outputs/poc_natural/cifar_sr_main/results.json [R-LOCAL] (100/124/800)
% input: outputs/poc_natural/cifar_duo_smoke/results.json [R-LOCAL] (duo 224)
% MRI NFE: protocol constants (T=20, q=2, R_a=2, N=4; Eq. eq:nfe), matching
% the pilot JSONs outputs/_ng_smoke_R4|R8.
% [length-edit] shell stripped: lives inside a paired table* in experiments.tex
  \centering\small
  \setlength{\tabcolsep}{3pt}
  \captionof{table}{Denoiser calls (\nfe) per reconstructed image for each
  acquisition-capable method, all measured in our pipeline\RLOCAL{}
  (Eq.~\ref{eq:nfe}; $NT$ for particle ensembles, $2T{+}2qR_a$ for
  \pvas-duo). The last column is the ratio to the passive
  single-trajectory baseline of the same domain.}
  \label{tab:nfe_cost}
  \begin{tabular}{llcc}
    \toprule
    Method & Uncertainty source & \nfe/img & vs.\ passive \\
    \midrule
    \multicolumn{4}{l}{\emph{CIFAR-10 block acquisition ($T{=}100$)}} \\
    Passive (static mask) & none & 100 & $\times$1.0 \\
    \pvas{} (ours) & 1-traj.\ Tweedie \textsc{pv} & 124 & $\times$1.24 \\
    \pvas-duo (ours) & \textsc{pv} $+$ disagreement & 224 & $\times$2.24 \\
    Ensemble-$8$ selection & $8$ posterior trajectories & 800 & $\times$8.0 \\
    \midrule
    \multicolumn{4}{l}{\emph{fastMRI multi-coil line acquisition ($T{=}20$)}} \\
    \pigdm{} (passive) & none & 20 & $\times$1.0 \\
    Static, matched budget & none & 20 & $\times$1.0 \\
    \pvas{} (ours) & 1-traj.\ Tweedie \textsc{pv} & 28 & $\times$1.4 \\
    Ensemble-$4$ selection & $4$ \pigdm{} particles & 80 & $\times$4.0 \\
    \bottomrule
  \end{tabular}
      % -> tab:nfe_cost
  \end{minipage}
\end{table*}

\paragraph{Results.} On the $n{=}128$ production run\RLOCAL{}
(Table~\ref{tab:sr_main}), \emph{pv} gains $+1.63$\,dB paired \psnrm{}
over \emph{static} (Wilcoxon $p{=}1.0{\times}10^{-21}$), $+1.45$\,dB
over \emph{random} ($p{=}7.1{\times}10^{-18}$), and $+2.05$\,dB over
\emph{raster} ($p{=}6.9{\times}10^{-22}$) --- at $1.24\times$ the \nfe,
recovering roughly half the gap to the retrospective \emph{oracle}. The
$8\times$-cost \emph{ensemble} remains stronger, but two thirds of its
margin persists under \emph{ens\_random}: it stems from particle-mean
averaging, not better selection. The budget sweep
($L\in\{6,12,18,24\}$, Fig.~\ref{fig:budget}; tabular form in the
supplement) is monotone: the paired margin of \emph{pv} over
\emph{static} grows from $+0.85$\,dB at $L{=}6$ to $+1.63$, $+2.59$,
and $+3.86$\,dB at $L{=}12,18,24$ (all
$p<4{\times}10^{-10}$)\RLOCAL{} --- tracking the
\emph{oracle}$-$\emph{static} headroom, which widens from $+1.8$ to
$+6.2$\,dB over the same range; we do not extrapolate beyond it. The
same protocol at $256^2$ scale is Sec.~\ref{sec:exp_celeba}.

\paragraph{Gate ablation.} The variance map can also replace the scalar
gain in data consistency (Eq.~\ref{eq:gate}); Table~\ref{tab:ablation_gate}
prices this byproduct on the SR task with noisy measurements
($\sigma_y{=}0.1$, $n{=}64$)\RLOCAL. Selection-only \pvas{} gains
$+1.64$\,dB over \emph{static} ($p{=}4.7{\times}10^{-12}$). Adding the
per-step gate on top of selection is \psnrm-neutral
($\Delta{=}{-}0.02$\,dB, $p{=}0.287$) at $2.2\times$ the cost. The
map's measured value is in \emph{selection}; the gate stays off
everywhere else in the paper.

\paragraph{Is one trajectory enough? (calibration).} Does a
single-trajectory Tweedie estimate rank candidate atoms as well as
cross-trajectory disagreement? We freeze a
shared partial trajectory at $50\%$ of the schedule ($50$ base \nfe)
and score all unobserved blocks with \pvas{} ($q\in\{1,2,4,8,16\}$
probes; $2q$ extra \nfe), ensemble variance from $N\in\{4,8,16\}$
trajectories re-run to the same point ($50N$ extra \nfe), and the
oracle; Fig.~\ref{fig:calibration} plots per-image
Spearman correlation with the oracle and top-$12$ overlap against extra
\nfe{} ($n{=}128$)\RLOCAL. In the SR regime the single-trajectory
signal is real and probe-limited: Spearman rises from $0.261$
($q{=}1$, $2$ extra \nfe) to $0.506$ ($q{=}16$, $32$ extra \nfe),
against $0.546$--$0.612$ for $N{=}4$--$16$ ensembles at $200$--$800$
extra \nfe; on top-$12$ overlap, $q{=}4$ reaches $0.357$ where the
$N{=}4$ ensemble reaches $0.412$ at $25\times$ the compute. In the
from-scratch inpainting regime the same signal is flat at Spearman
$0.03$--$0.08$ across all probe counts, with chance-level overlap
($0.18$--$0.20$ vs.\ chance ${\approx}0.21$), while the $N{=}16$
ensemble reaches $0.283$/$0.307$ --- the boundary condition analyzed
in Sec.~\ref{sec:exp_inpaint}.

\subsection{Boundary condition: from-scratch inpainting, and the duo
rescue}\label{sec:exp_inpaint}

Replacing the always-observed LR preview with a sparse random
initialization ($10\%$ of blocks, no preview) yields from-scratch
inpainting --- and single-trajectory \pvas{} stops working. We report
this negative result in full in Table~\ref{tab:inpaint} (same protocol,
$n{=}128$)\RLOCAL: \emph{pv} falls $1.54$\,dB below \emph{static} and
is indistinguishable from \emph{random} mid-trajectory acquisition
(paired $\Delta{=}{-}0.02$\,dB, $p{=}0.85$), while the retrospective
\emph{oracle} still clears \emph{static}. Two mechanisms are at work.
First, an uninformative signal exposes a pure \emph{timing} penalty:
up-front blocks condition the whole trajectory, blocks arriving at
$45$--$75\%$ of the schedule cannot undo an early commitment, and both
active policies land ${\approx}1.5$\,dB below \emph{static}. Second,
the ensemble's win here is averaging, not selection: \emph{ensemble}
itself falls \emph{below} \emph{ens\_random} (Table~\ref{tab:inpaint}),
and Fig.~\ref{fig:calibration} shows PV ranking at chance level in this
regime for every probe count. The failure is exactly the one
Proposition~1 predicts. With no anchor, the trajectory commits early to
one posterior mode and hallucinates coherent content; the denoiser
Jacobian then measures within-mode curvature around the committed
sample --- small precisely where the sample is confidently wrong ---
while the cross-mode term of Eq.~\ref{eq:ltv} carries the signal.

\paragraph{\pvas-duo restores the missing term.} The decomposition also
says what the minimal fix costs: one more trajectory. \pvas-duo
(Eq.~\ref{eq:duo}) adds to the within-mode score the squared
disagreement of two synchronized trajectories --- an unbiased
two-sample estimate of the cross-mode variance --- at
$2T{+}2qR_a{=}224$ \nfe{} (${\approx}2.2\times$ passive). In an
$n{=}16$ pilot on this task\RLOCAL{}, duo reaches $14.76$ \psnrm{}
where \emph{static} reaches $13.83$ ($+0.94$\,dB) and
single-trajectory \emph{pv} collapses to $12.70$; the $N{=}8$ ensemble
remains strongest at $16.60$ for $800$ \nfe. The ladder of
Sec.~\ref{sec:method_when} is visible end to end: one trajectory
suffices under strong conditioning (SR, MRI), two restore active
selection under weak conditioning, and a full ensemble is what severe
multi-modality actually costs. We believe this is why
AdaSense~\cite{elata2024adasense}, ADS~\cite{nolan2025ads}, and
diffusion-BED~\cite{iollo2025bedmri} pay for ensembles --- and that for
conditioned acquisition problems they overpay.

\subsection{Guided super-resolution on CelebA-HQ
$256^2$}\label{sec:exp_celeba}

The same protocol scales to $256{\times}256$ faces
\cite{karras2018progressive} under a pretrained unconditional DDPM
\cite{ho2020ddpm}. The $8\times$ box-downsampled $32^2$ preview is
always observed; atoms are $8{\times}8$ blocks ($32{\times}32$ grid,
$1024$ atoms), and the budget buys $L{=}96$ of them ($9.4\%$);
$T{=}100$, $R_a{=}3$, $q{=}4$, $n{=}64$ fixed test images.
Table~\ref{tab:celeba} reports the policy set of
Table~\ref{tab:sr_main} plus \pvas-duo, with the particle ensemble at
$N{=}4$ ($400$ \nfe); qualitative panels are in the supplement
(Fig.~S1).

% Hand-written shell for the CelebA-HQ 256 quantitative table (tab:celeba).
% Metric cells merge from outputs/poc_natural/celeba_quant/results.json
% (poc_natural/quantitative_celeba.py) when that file lands; NFE column is
% protocol-derived (T=100, q=4, R_a=3, N=4; Eq. eq:nfe / 2T+2qR_a for duo).
% input: outputs/poc_natural/celeba_quant/results.json [MISSING -> \PENDING]
% [length-edit] hand-maintained: gen_tables.py must write .new for merge.
\begin{table}[t]
  \centering\small
  \setlength{\tabcolsep}{4pt}
  \caption{Guided super-resolution on CelebA-HQ $256^2$ at matched
  measurement budget: the $8\times$ box-downsampled $32^2$ preview is
  always observed; the budget buys $L{=}96$ of $1024$ atoms
  ($8{\times}8$-px blocks, $9.4\%$); $T{=}100$ steps, $n{=}64$ images.
  Oracle and ensemble rows are references. ``\PENDING'' marks production
  cells whose result files were not yet written when this text was
  frozen (Sec.~\ref{sec:exp}).}
  \label{tab:celeba}
  \begin{tabular}{lccc}
    \toprule
    Acquisition policy & \psnrm{} (dB) & \ssimm{} & \nfe/image \\
    \midrule
    Static mask (random, up-front) & \PENDING & \PENDING & 100 \\
    Random (mid-trajectory) & \PENDING & \PENDING & 100 \\
    \pvas{} (ours) & \PENDING & \PENDING & 124 \\
    \pvas-duo (ours) & \PENDING & \PENDING & 224 \\
    Oracle error (upper bound) & \PENDING & \PENDING & 100 \\
    Ensemble-$4$ variance (ref.) & \PENDING & \PENDING & 400 \\
    \midrule
    \multicolumn{4}{l}{\emph{Paired $\Delta$\psnrm{} (\pvas{} $-$ baseline), dB}} \\
    $\Delta$ vs.\ static & \PENDING & \PENDING & \\
    $\Delta$ vs.\ random & \PENDING & \PENDING & \\
    \bottomrule
  \end{tabular}
\end{table}
  % -> tab:celeba

\subsection{Multi-coil MRI line acquisition}\label{sec:exp_mri}

\paragraph{Setup.} We use $10$ multi-coil brain volumes (AXT2) from
the fastMRI validation set \cite{zbontar2018fastmri}, $3$ central
slices each ($n{=}30$), coil images center-cropped to $384{\times}320$;
1D Cartesian column undersampling at $R\in\{4,8\}$ with an $8\%$ ACS
band (at $R{=}4$ the base mask has $85$ columns\RLOCAL). Per-coil noise
variances are estimated from $k$-space corners, and complex white noise
is added retrospectively with per-slice deterministic seeds, identical
across methods; coil sensitivities, estimated once from the
low-frequency ACS of the noisy data \cite{uecker2014espirit}, are
shared by all methods. The prior is a single unconditional EDM denoiser
\cite{karras2022edm} trained on fully sampled fastMRI brain
reconstructions and frozen; sampling uses a $T{=}20$-step EDM
$\sigma$-schedule from $x_T=\sigma_{\max}\,\varepsilon$, truncated at
$\sigma_{\max}{=}10$ (headline) with the EDM-default
$\sigma_{\max}{=}80$ regime in the supplement. \pvas{} runs
\emph{selection-only} (Sec.~\ref{sec:method_inst}): data consistency
uses the scalar \pigdm{} gain everywhere, and the estimator
(Eq.~\ref{eq:pvk}, $q{=}2$ probes) fires only at the $R_a{=}2$
acquisition rounds, at $50\%$ and $85\%$ of the schedule. Active runs
acquire $L\in\{8,16\}$ columns; the cost is $\nfe{=}T+2qR_a{=}28$
against $20$ for a passive solve ($1.4\times$).

\paragraph{Methods.} \emph{pigdm\_mc}: multi-coil \pigdm{}
\cite{song2023pigdm} on the base $R$-mask --- the passive solver
\pvas{} builds on. \emph{static\_match}: \pigdm{} on a static random
mask at the full matched budget $\lvert M_R\rvert{+}L$ including the
ACS band. \emph{random\_adaptive}: identical schedule and budget with
uniform-random columns and no probes, isolating variance-\emph{guided}
selection from generic mid-trajectory adaptation. \emph{pv\_active}
(\pvas): PV-scored acquisition. \emph{oracle\_active}: columns scored
by the residual against the held-out fully sampled $k$-space --- the
ceiling for any residual-driven score, not deployable.
\emph{ens\_active}: the faithful ensemble baseline of
Sec.~\ref{sec:method_active} --- $N{=}4$ parallel \pigdm{} particles
sharing the measurement state, columns scored by the cross-particle
variance of the particles' predicted coil $k$-space (the
measurement-space disagreement signal of ADS and AdaSense
\cite{nolan2025ads,elata2024adasense}), reconstruction the particle
mean, $\nfe{=}NT{=}80$. \emph{ens\_random\_active}: the same $N{=}4$
particles with random selection (the averaging control). Metrics:
\psnrm/\ssimm{} on the central $320^2$ magnitude crop, SENSE-referenced
(the SENSE combination of clean, fully sampled coil data under the
\emph{same} estimated maps, removing map-estimation error common to all
methods); RSS-referenced values are in the released result files.

% Hand-restructured for the selection-only (gate-off) production protocol
% (scripts/reconstruct_mc.py, active_gate=off, q=2, R_a=2). Metric cells
% merge from outputs/PVAS_ng/R{4,8}_a{8,16}/results.json when those land
% (point gen_tables.py mri_table at PVAS_ng); Cols and NFE are protocol
% constants (base masks: 85 cols at R=4, 43 at R=8, measured in the pilot
% JSONs outputs/_ng_smoke_R4|R8; NFE = T + 2qR_a = 28, ensembles N*T = 80).
% input: outputs/PVAS_ng/R4_a8/results.json  [MISSING -> \PENDING]
% input: outputs/PVAS_ng/R4_a16/results.json [MISSING -> \PENDING]
% input: outputs/PVAS_ng/R8_a16/results.json [MISSING -> \PENDING]
% The former gate-on grid (outputs/PVAS_local) is now one ablation sentence
% in Sec. exp_mri, no longer this table.
% [length-edit] compact styling + R8_a8 group lives in the supplement (tab:mri_r8)
\begin{table}[t]
  \centering\footnotesize
  \renewcommand{\arraystretch}{0.95}
  \setlength{\tabcolsep}{3pt}
  \caption{Multi-coil fastMRI brain k-space line acquisition,
  selection-only protocol ($\sigma_{\max}{=}10$, $T{=}20$ EDM steps,
  $q{=}2$ probes at $R_a{=}2$ rounds; $n{=}30$ slices per configuration).
  All methods share the static baseline's total column budget;
  SENSE-referenced \psnrm/\ssimm; \nfe{} per Eq.~\ref{eq:nfe} ($NT$ for
  particle ensembles). $\Delta$: paired per-slice differences, one-sided
  Wilcoxon $p$. ``\PENDING'' marks production cells whose result files
  were not yet written when this text was frozen; the two-slice pilot of
  the identical protocol is quoted in Sec.~\ref{sec:exp_mri}. $R{=}8$,
  $L{=}8$ and $\sigma_{\max}{=}80$: supplement.}
  \label{tab:mri_main}
  \begin{tabular}{lcccc}
    \toprule
    Method & Cols & \psnrm{} (dB) & \ssimm{} & \nfe/slice \\
    \midrule
    \multicolumn{5}{l}{\emph{$R{=}4$, $L{=}8$ active lines}} \\
    \pigdm{} (passive) & 85 & \PENDING & \PENDING & 20 \\
    Static, matched budget & 93 & \PENDING & \PENDING & 20 \\
    Random adaptive & 93 & \PENDING & \PENDING & 20 \\
    \pvas{} (ours) & 93 & \PENDING & \PENDING & 28 \\
    Ensemble-$4$ active (ref.) & 93 & \PENDING & \PENDING & 80 \\
    Ensemble-$4$, random select (ref.) & 93 & \PENDING & \PENDING & 80 \\
    Oracle selection (bound) & 93 & \PENDING & \PENDING & 20 \\
    $\Delta$(\pvas{} $-$ static match) & ($n{=}30$) & \PENDING & & \\
    $\Delta$(\pvas{} $-$ random adaptive) & ($n{=}30$) & \PENDING & & \\
    \midrule
    \multicolumn{5}{l}{\emph{$R{=}4$, $L{=}16$ active lines}} \\
    \pigdm{} (passive) & 85 & \PENDING & \PENDING & 20 \\
    Static, matched budget & 101 & \PENDING & \PENDING & 20 \\
    Random adaptive & 101 & \PENDING & \PENDING & 20 \\
    \pvas{} (ours) & 101 & \PENDING & \PENDING & 28 \\
    Ensemble-$4$ active (ref.) & 101 & \PENDING & \PENDING & 80 \\
    Ensemble-$4$, random select (ref.) & 101 & \PENDING & \PENDING & 80 \\
    Oracle selection (bound) & 101 & \PENDING & \PENDING & 20 \\
    $\Delta$(\pvas{} $-$ static match) & ($n{=}30$) & \PENDING & & \\
    $\Delta$(\pvas{} $-$ random adaptive) & ($n{=}30$) & \PENDING & & \\
    \midrule
    \multicolumn{5}{l}{\emph{$R{=}8$, $L{=}16$ active lines}} \\
    \pigdm{} (passive) & 43 & \PENDING & \PENDING & 20 \\
    Static, matched budget & 59 & \PENDING & \PENDING & 20 \\
    Random adaptive & 59 & \PENDING & \PENDING & 20 \\
    \pvas{} (ours) & 59 & \PENDING & \PENDING & 28 \\
    Ensemble-$4$ active (ref.) & 59 & \PENDING & \PENDING & 80 \\
    Ensemble-$4$, random select (ref.) & 59 & \PENDING & \PENDING & 80 \\
    Oracle selection (bound) & 59 & \PENDING & \PENDING & 20 \\
    $\Delta$(\pvas{} $-$ static match) & ($n{=}30$) & \PENDING & & \\
    $\Delta$(\pvas{} $-$ random adaptive) & ($n{=}30$) & \PENDING & & \\
    \bottomrule
  \end{tabular}
\end{table}
  % -> tab:mri_main

\paragraph{Results.} Table~\ref{tab:mri_main} carries the production
grid ($R\in\{4,8\}$, $L\in\{8,16\}$, $n{=}30$ slices per
configuration); it is the evidence of record for this section. A
two-slice pilot of the identical protocol\RLOCAL{} ($n{=}2$; result
files ship with the code) shows the expected ordering. At $R{=}8$,
$L{=}16$: \pvas{} reaches $30.27$ SENSE-\psnrm{} against $28.33$ for
passive \pigdm{} ($+1.94$\,dB), $28.90$ for \emph{static\_match}
($+1.38$), and $29.06$ for \emph{random\_adaptive} --- matching the
residual oracle ($30.30$) and edging the $N{=}4$ particle ensemble
($30.18$) at $28$ vs.\ $80$ \nfe. At $R{=}4$, $L{=}8$ the margins
shrink and the ensemble leads narrowly: \pvas{} $32.74$ vs.\ $31.94$
\pigdm{} ($+0.80$), $31.55$ static ($+1.19$), $32.34$ random; oracle
$32.93$; ensemble $32.91$. At $n{=}2$ these are smoke checks, not
evidence; we attach no significance to them.

\paragraph{Gate-on ablation.} An earlier variant of this grid refreshed
the variance map every step to drive the Kalman gain of
Eq.~\ref{eq:gate} ($100$ \nfe). Selection gains persisted there
($n{=}30$ per configuration: $+0.37$\,dB over static at $R{=}4$,
$L{=}8$; $+0.11$\,dB n.s.\ at $R{=}4$, $L{=}16$; $+1.16$\,dB at
$R{=}8$, $L{=}16$)\RLOCAL, but the gate itself sat
\psnrm-neutral-to-negative against the scalar \pigdm{} gain on the same
masks --- consistent with the gate ablation of Sec.~\ref{sec:exp_sr}
--- so production runs selection-only at $28$ \nfe.

\subsection{Cost analysis}\label{sec:cost}

Table~\ref{tab:nfe_cost} prices acquisition-grade uncertainty in
denoiser calls (Eq.~\ref{eq:nfe}), all measured in our
pipeline\RLOCAL. \pvas{} spends $124$ \nfe{} at the CIFAR configuration
($1.24\times$ passive) and $28$ for MRI ($1.4\times$); \pvas-duo
doubles the trajectory ($224$, ${\approx}2.2\times$); ensemble-based
selection multiplies the \emph{entire} trajectory by $N$ --- $800$
\nfe{} for the $N{=}8$ CIFAR reference, $80$ for the $N{=}4$ MRI
reference, run head-to-head at matched budget rather than transplanted
from a different prior and protocol. For context, the published systems
prescribe larger ensembles still: $N_p{=}16$ parallel particles in
ADS\CITED~\cite{nolan2025ads}, up to $256$ complete posterior-sampling
runs per image in AdaSense\CITED~\cite{elata2024adasense}, and an
ensemble of trajectories per design step in
diffusion-BED\CITED~\cite{iollo2025bedmri}. The honest converse
(Sec.~\ref{sec:exp_inpaint}, Fig.~\ref{fig:calibration}): ensembles buy
robustness to posterior multi-modality that no probe count recovers
from one trajectory, and duo prices the first rung of that ladder.
